\section{How much splitting information can congruences recover?}
\label{sec:rigidity}

A common route from Dirichlet GRH to arithmetic over number fields is to
express a Dedekind zeta function as a product of Dirichlet $L$-functions.
There is an exact obstruction to applying this route to a nonabelian normal
closure. It can be measured either by the frequency of incorrect splitting
predictions or by the mean squared error in the first Euler coefficients.
The formulas below follow from Chebotarev and Kronecker--Weber. We claim
no priority for them. For the classical splitting problem, see
\href{https://arxiv.org/abs/2208.11381}{Rosen, Takeyama, Tasaka and Yamamoto};
for the input theorems see
\href{https://www.jmilne.org/math/CourseNotes/cft.html}{Milne, \emph{Class Field Theory}}.

\paragraph{Best possible congruence test.}
Let $L/\mathbb Q$ be finite Galois with group $G$, let $D=[G,G]$, and let
$C\subseteq G$ be a union of conjugacy classes. One wants to decide whether
$\operatorname{Frob}_p\in C$ from a fixed residue $p\bmod m$. Allow
\emph{any} fixed modulus $m$ and \emph{any} choice of acceptable residue
classes. The minimum natural density, relative to all primes, of incorrect
decisions is exactly
\[
 \boxed{\displaystyle
  \frac{1}{|G|}\sum_{T\in G/D}\min\{\,|T\cap C|,\ |T\setminus C|\,\}.}
\]
The minimum is attained. Ramified primes and primes dividing the modulus
form a finite exceptional set.

\emph{Proof.} Put $M=\mathbb Q(\zeta_m)$ and $E=L\cap M$.
The Galois group of $LM$ consists of the compatible pairs in
$G\times\operatorname{Gal}(M/\mathbb Q)$ with equal restriction to $E$.
By Chebotarev, prime densities for any conjugacy-invariant test on these
pairs are uniform element counts. Fixing the second coordinate leaves the
first uniformly distributed, for such tests, in a coset of
$\operatorname{Gal}(L/E)$. Since $E$ is abelian, this coset is a union
of $D$-cosets. Even if one learns the individual $D$-coset, the majority
decision on it still makes at least the displayed number of errors.
Kronecker--Weber supplies an $m$ for which $L^D\subseteq M$. Then
$E=L^D$ and the residue reveals exactly the required coset. Choosing its
majority value attains the bound.\hfill$\square$

For example, let $K$ be any number field and take $L$ to be its normal closure.
A prime splits completely in $K$ precisely when its Frobenius acts trivially
on all embeddings of $K$. This action is faithful, so $C=\{1\}$. The minimum
error density equals $0$ for abelian $G$ and $1/|G|$ for nonabelian $G$.
In the second case the rule declaring that no prime splits completely attains
the minimum. In particular, even the best rule for complete splitting in
$\mathbb Q(\sqrt[3]{2})$ makes errors on a set of density at least $1/6$.
The splitting field has group $S_3$, since $X^3-2$ is irreducible by Eisenstein
and has nonsquare discriminant $-108$.

\paragraph{Best approximation to the Euler coefficient.}
Write $H=\operatorname{Gal}(L/K)$ and $F=L^{HD}$, the largest subfield of $K$
abelian over $\mathbb Q$. For an unramified prime $p$, write $a_K(p)$ for the
number of degree-one primes of $K$ above $p$. The optimum over all
complex-valued periodic functions $b$ is
\[
 \boxed{\displaystyle
 \min_b\lim_{x\to\infty}\frac{1}{\pi(x)}\sum_{p\le x}
          |a_K(p)-b(p)|^2
      = |H\backslash G/H|-[F:\mathbb Q].}
\]
An optimal $b$ equals $a_F(p)$ away from the finite exceptional set.

\emph{Proof.} Let $u(g)$ count the fixed points of $g$ on $G/H$. At every
unramified $p$ this is $a_K(p)$ for $g=\operatorname{Frob}_p$. By the
same compatible-pair argument, the sharpest available conditional mean
is the average of $u$ on each $D$-coset, attained at a modulus containing
$L^D$ in its cyclotomic field. Indeed a squared error on a finite set is
minimized by its mean, and a coarser partition cannot reduce the minimum.
The coset average is the character of the $D$-invariants of the permutation
representation on $G/H$, because averaging the $D$-action projects onto
this subspace. A basis is indexed by $D$-orbits, which form the
$G$-set $G/HD$. Its character is precisely $a_F$ on Frobenii. This
character is $[G:HD]$ on $HD$ and zero elsewhere, since $HD$ is normal.
Its squared norm for normalized counting measure on $G$ equals $[G:HD]$.
The squared norm of $u$ counts $G$-orbits on $(G/H)^2$ by Burnside's lemma.
Those orbits correspond to the $H$-double cosets in $G$. Subtracting
the norm of the orthogonal projection gives the formula.\hfill$\square$

For $\mathbb Q(\sqrt[3]{2})$ the best approximation is the constant $1$.
The values $3,1,0$ occur with densities $1/6,1/2,1/3$, giving squared
error $4/6+1/3=1$. The sign of Frobenius, which gives the quadratic
resolvent, distinguishes odd from even permutations but does not improve
this squared error. The file \texttt{checks/rigidity.py} checks these
finite group calculations and the corresponding ones for cyclic quartic,
dihedral quartic and $S_4$ quartic permutation actions.

\paragraph{Consequence for direct transfer.}
Suppose $\zeta_K$ agrees in $\Re s>1$ with a finite product of integer powers
of ordinary Dirichlet $L$-functions up to Euler factors at finitely many primes.
By uniqueness of Dirichlet coefficients, its coefficient at almost every $p$
equals a fixed linear combination of values $\chi(p)$. This is periodic in $p$,
so the condition $a_K(p)=[K:\mathbb Q]$ would predict complete splitting
with zero error density. The first formula forces $G$ to be abelian.
Conversely, abelian fields over $\mathbb Q$ have precisely such a
factorization via their primitive Dirichlet characters. Thus any attempt
to establish GRH for the zeta function of a nonabelian normal closure,
or of its nonnormal generating subfield, needs an argument beyond a
finite Dirichlet product or quotient. This obstruction concerns that
particular proof strategy. Zeros of distinct analytic functions can still
be related by other means.

\paragraph{A one-sided criterion for the missing zeros.}
There is a concrete use for the optimal predictor $a_F$. Put
\[
 U_K(x)=\sum_{p\le x,\ p\text{ unramified in }L}
       (a_K(p)-a_F(p))\log p\,(1-p/x).
\]
With the repository's Dirichlet GRH as input, the following are equivalent:
\[
 \mathrm{RH}(\zeta_K)
 \quad\Longleftrightarrow\quad
 \exists C_K\ \ \forall n\in\mathbb Z,\ n\ge2:\quad
 U_K(n)\ge-C_K\sqrt n .
\]
This works for nonnormal $K$ as well. In the primitive $A_5$ sextic
case, $F=\mathbb Q$ since $A_5$ is perfect. The weights $a_K(p)-1$
are respectively $5,1,-1,0$ at Frobenius elements of orders
$1,2,3,5$. Thus the missing inequality is a single lower bound for a
sum over rational primes with four explicit weights.

\emph{Proof.} Write $V_J(x)=\sum_{N\mathfrak p^r\le x}\log N\mathfrak p
(1-N\mathfrak p^r/x)$, summing over prime ideals of a field $J$ and
integers $r\ge1$. Then $V_J=x^{-1}\int_1^x\psi_J(u)\,du$. Under field RH
the integrated explicit formula yields $V_J=x/2+O_J(\sqrt x)$: the
nontrivial zero sum after division by $x$ is bounded by
$\sqrt{x}\sum_\rho |\rho(\rho+1)|^{-1}$, a convergent sum.
See \href{https://arxiv.org/abs/1312.4465}{Greni\'e--Molteni,
formula (1.3a)}. Prime powers of rational exponent at least two
make a contribution $O_J(\sqrt x)$, by bounding their log coefficients
by $[J:\mathbb Q]\log p$ and using the ordinary Chebyshev bound,
first for squares and then for exponents at least three. The finitely many
ramified primes cost $O_J(\log x)$. Thus $U_K=V_K-V_F+O_L(\sqrt x)$.
Here $F$ satisfies RH by abelian factorization. This proves necessity.

For sufficiency extend the integer lower bound to all real $x\ge2$.
Between primes $|U_K'(x)|\ll_K x^{-2}\sum_{p\le x}p\log p\ll_K 1$,
and $U_K$ is continuous. Write $W=V_K-V_F$. For $\Re s>1$,
\[
 \int_1^\infty W(x)x^{-s-1}\,dx
 =\frac{-\zeta'_K(s)/\zeta_K(s)+\zeta'_F(s)/\zeta_F(s)}{s(s+1)} .
\]
The pole terms at $1$ cancel. Suppose $\zeta_K$ has a zero $\rho$ with
$\Re\rho>1/2$. It gives a genuine pole since $\zeta_F$ is zero-free there.
Add $C\sqrt x$ to $W$ to make it nonnegative, increasing $C$ on a compact
interval if necessary. This adds $C/(s-1/2)$ to the transform and preserves
the pole at $\rho$. The integral's real abscissa of convergence $b$
therefore satisfies $\Re\rho\le b\le1$. Landau's theorem forces a
singularity at the real point $b$: Taylor expansion from the right has
nonnegative moments of $\log x$ as coefficients in the direction of
decreasing $s$, so analyticity at $b$ would give convergence across $b$.
At $b=1$ the transform is regular. At any real zero $b\in(1/2,1)$ of
multiplicity $h$, its residue is $-h/[b(b+1)]<0$, inconsistent with the
nonnegative integral to the right. These exhaust the possible real
singularities. The functional equation gives the other half of the
critical strip.\hfill$\square$

The incoming reports apply a related sign argument to completely split
primes in a Galois field and to individual cubic examples. The conditional
mean here selects the abelian contribution for every field and the smoothing
removes the pole without a field-dependent centering constant. Proving the
displayed bound for the primitive sextic family still requires control of
the nonabelian weights.
